% Tipo de documento. En este caso es un art�culo, para folios A4, tama�o de la fuente 11pt y con p�gina separada para el t�tulo
\documentclass[a4paper,11pt,titlepage]{article}

% Carga de paquetes necesarios. OrdenesArticle es un paquete personalizado
\usepackage[spanish]{babel} 
\RequirePackage[T1]{fontenc}
\RequirePackage[ansinew]{inputenx} 
\usepackage[spanish,cap,cont,title,fancy]{OrdenesArticle}
\usepackage{lmodern}
\usepackage{array}
\usepackage{graphicx}
\usepackage{hyperref}
\usepackage{pifont}
\usepackage{listings}
\usepackage[usenames,dvipsnames]{color}
\usepackage{colortbl}
\usepackage{color}
\usepackage{ifthen}
\usepackage{longtable}
\hypersetup{bookmarksopen,bookmarksopenlevel=4,linktocpage,colorlinks,urlcolor=blue,citecolor=blue,
						linkcolor=blue,filecolor=blue,pdfnewwindow,
						pdftitle={PlanCoDE: Planificador Colaborativo para el Dise�o de Estrategias y Acciones de Emergencia},
						pdfauthor={Juan Andrada Romero, Jose Domingo L�pez L�pez},
						pdfsubject={Sistemas para la Colaboraci�n}}


% Macro para definir una lista personalizada 
\newenvironment{milista}%
{\begin{list}{\textbullet}%
{\settowidth{\labelwidth}{\textbullet} \setlength{\leftmargin}{\dimexpr\labelsep+\labelwidth+5pt}
\setlength{\itemsep}{\dimexpr 0.5ex plus 0.25ex minus 0.25ex}
\setlength{\parsep}{\itemsep}
\setlength{\partopsep}{\itemsep}
\addtolength{\topsep}{-7.5pt}
}}%
{\end{list}}

% Macro para insertar una imagen
%       Uso: \imagen{nombreFichero}{Factor escala}{Caption (leyenda)}{Label (identificador para referenciarla)}
% -------------------------------------------------------------------------------------------------------------
\def\imagen#1#2#3#4{
 \begin{figure}[h]
 \begin{center}
   \scalebox{#2}{\includegraphics{#1}}
 \caption {#3}
 \label{#4}
 \end{center}
 \end{figure}
}

% Macro para dar color a las filas de la tabla "`Operaciones"'
\def\colorFila{
	\rowcolor[rgb]{0.4,0.52,1.0}
}

% Definicion del "listings" para el lenguaje Java
\lstdefinelanguage{Java}
{
 morecomment = [l]{//}, 
 morecomment = [l]{///},
 morecomment = [s]{/*}{*/},
 morestring = [b]", 
 sensitive = true,
 morekeywords = {package, static, while, switch, break, line, void, String, Object, int, Integer, instanceof, else, if, for, private, return, new, public, class, import, int, boolean, true, false, extends, final, super, protected, abstract, this, do, float, double, null, try, catch, implements}
}

% Configuraci�n del c�digo incluido para el lenguaje Java
\lstset{
  language=Java,
  basicstyle=\footnotesize,
  backgroundcolor=\color{white},
  showspaces=false,
  showstringspaces=false,
  showtabs=false,
  frame=single,
  tabsize=2,
  captionpos=b,
  breaklines=true,
  breakatwhitespace=false,
  escapeinside={\%},
  keywordstyle = \color [rgb]{0,0,1},
  commentstyle = \color [rgb]{0.133,0.545,0.133},
  stringstyle = \color [rgb]{0.627,0.126,0.941}
}

\begin{document}

% En las p�ginas de portada e �ndices, no hay encabezado ni pie de p�gina
\pagestyle{empty} 

% Se incluye la portada
\begin{titlepage}
	\begin{center}
  	{\LARGE UNIVERSIDAD DE CASTILLA-LA MANCHA} \\
  	\bigskip
  	{\Large ESCUELA SUPERIOR DE INFORM�TICA} \\
  	\vspace{20mm}
  	\includegraphics[scale=0.45, keepaspectratio]{./images/esi_bw.png} \\
  	\vspace{20mm}
  	{\Huge \textbf{Almacenamiento y Recuperaci�n de la Informaci�n}}\\
  	\vspace{10mm}
  	{\LARGE \textsc{\textbf{M�dulo de importaci�n e indexaci�n de documentos}}}\\
  	\vspace{30mm}
  	\today
  	\vspace{40mm}
  	\begin{flushleft}
  		{\large Juan Andrada Romero}\\
  		\vspace{1mm}
  		{\large Jose Domingo L�pez L�pez}\\		
  	\end{flushleft}
	\end{center}
\end{titlepage}
\clearpage

% Texto del reverso de la portada
%%\mbox{}
%%\vspace{18cm}
%%\begin{small}
% Se ajusta la separaci�n entre p�rrafos
%%\parskip=10pt 

%%\copyright~ 2008/2009 Juan Andrada Romero. Universidad de Castilla La Mancha, Escuela Superior de Inform�tica de Ciudad Real.

%%Se permite la modificaci�n, copia y distribuci�n de este documento, seg�n la licencia de documentaci�n GNU (\url{http://www.gnu.org}).

%%Este documento fue compuesto con \LaTeX{}. Im�genes generadas con Power Point y Gimp.
%%\end{small}

%%\newpage

% En las p�ginas de �ndices y prefacio, se utiliza numeraci�n romana
\pagenumbering{Roman}

% Se crea el �ndice
\tableofcontents
% Se pasa p�gina y se a�ade el �ndice de figuras al �ndice principal
%\clearpage\phantomsection
%\addcontentsline{toc}{section}{\listfigurename}
% Se crea el �ndice de figuras
%\listoffigures
% Si se quiere crear un �ndice de tablas se pondr�a: \listoftables

\newpage

% Se ajusta la separaci�n entre p�rrafos
\parskip=10pt

% Se a�ade el prefacio al �ndice
%\clearpage\phantomsection
%\addcontentsline{toc}{section}{Prefacio}

% Comienza el contenido del documento. Se utilizan n�meros ar�bigos y el encabezado y pie de p�gina personalizado
\pagenumbering{arabic}
\pagestyle{fancy}

\section{Introducci�n}


\section{�Qu� es PlanCoDE?}

\subsection{Diagrama de casos de uso}

\section{Clasificaci�n} \label{clasificacion}

En los siguientes apartados, se realizar� una clasificaci�n de la herramienta PlanCode, seg�n los diferentes autores que aparecen en \cite{groupware}.

\subsection{Dimensiones espacio-temporal}

Atendiendo a la clasificaci�n en tiempo y espacio propuesta por \textbf{Johansen}, PlanCode puede clasificarse en la categor�a \textit{diferente lugar, mismo tiempo}, es decir, es una herramienta de interacci�n s�ncrona distribuida (ver Figura \ref{fig:johansen}). Esto es as� porque los diferentes usuarios deben estar conectados en la misma sesi�n y al mismo tiempo para poder recibir los mensajes del chat, la imagen del mapa cargada por alguno de ellos y los trazos que se van realizando.

\imagen{./imagenes/johansen.png}{0.3}{Taxonom�a en tiempo y espacio (Johansen; 1991)}{fig:johansen}

Por otra parte, atendiendo a la clasificaci�n espacio-temporal propuesta por \textbf{Grudin}, la herramienta se puede clasificar tambi�n en la categor�a \textit{diferente lugar, mismo tiempo}, por las mismas razones que en el caso anterior. 


\subsection{Dominio de aplicaci�n}

Siguiendo la clasificaci�n del nivel de aplicaciones de \textbf{Ellis}, PlanCode cumple con las caracter�sticas de los siguientes sistemas:
\begin{enumerate}[A.]
	\item \textbf{Sistema de mensajes}: PlanCode se puede considerar un sistema de mensajes porque permite el intercambio de mensajes de diferente naturaleza entre los usuarios, como es el intercambio de mensajes de texto a trav�s del chat, im�genes cuando se carga un mapa y el intercambio de trazos libres.
\end{enumerate}

\subsection{Modelo 3C}

Si atendemos al objetivo principal de PlanCode (el dise�o de rutas sobre un mapa en situaciones de emergencias), dicha herramienta se puede clasificar como una herramienta de colaboraci�n, pues, como ya se ha comentado, los trazos pueden ser dibujados por cualquier usuario, apoyando al resto de trazos que ya est�n dibujados. Del mismo modo, cualquier usuario podr�a modificar la imagen del mapa que se est� mostrando si lo cree conveniente. \\
\indent Pero, como PlanCode cuenta con un chat, se puede considerar tambi�n una herramienta de comunicaci�n y coordinaci�n.



\section{Requisitos del sistema}

Para poder utilizar \textbf{FAQ Manager} ser� necesario disponer de un servidor de p�ginas Web (por ejemplo, Apache HTTP Server \cite{apache-homepage}) que pueda procesar c�digo PHP \cite{php-homepage}. Adem�s, ser� necesario un sistema gestor de bases de datos como MySQL \cite{mysql-homepage}. No obstante, para simplificar el proceso de instaci�n y configuraci�n de todo este software, se recomienda el uso de WAMP \cite{wamp-homepage}.

Una vez instalado y configurado el software necesario, deber� editar el fichero \texttt{config.php} que se incluye en la distribuci�n de \textit{FAQ Manager}. En este fichero, de estructura similiar a la de Figura \ref{fig:config}, deber� cumplimentar debidamente el dominio sobre el que se aloja el sistema gestor de bases de datos, el nombre de la base de datos, el nombre de un usuario de la base de datos con los permisos b�sicos para realizar operaciones \textit{CRUD} (Create, Read, Update, Delete) y su contrase�a.

\imagen{./imagenes/config.png}{0.7}{Fichero \texttt{config.php}}{fig:config}

\section{Prototipo de interfaz gr�fica de usuario}

% Secci�n de ap�ndices
% Se a�ade el titulo de "Ap�ndices" al �ndice
%\addcontentsline{toc}{section}{Ap�ndices}

%\input{Apendices.tex}

% A�adimos la bibliografia al �ndice
\clearpage\phantomsection
\addcontentsline{toc}{section}{Referencias}

% Bibliograf�a. En este caso se usa BibTeX
%\bibliographystyle{plain}
%\pagestyle{plain} 
%\bibliography{Bibliografia}

\end{document}
